% Compile with pdfLaTeX or upload to Overleaf.
\documentclass[11pt,a4paper]{article}
\usepackage[utf8]{inputenc}
\usepackage[T1]{fontenc}
\usepackage{lmodern}
\usepackage{geometry}
\usepackage{amsmath,amssymb,amsthm,mathtools}
\usepackage{listings}
\usepackage{hyperref}
\geometry{margin=1in}
\hypersetup{colorlinks=true,linkcolor=black,citecolor=black,urlcolor=blue}
\lstset{basicstyle=\ttfamily\small,breaklines=true,frame=single,columns=fullflexible}

\title{Mathematical Foundations of AES and\\
Finite-Field-Based Master-Key Preprocessing}
\author{}
\date{}

\begin{document}
\maketitle

\begin{abstract}
This report studies several algebraic structures used in AES, beginning with
the construction of $\mathrm{GF}(256)$ and polynomial arithmetic over
$\mathrm{GF}(2)$. It then explains multiplicative inversion, construction of
the AES S-box, AES-128 key expansion, and the standard round transformations.
Finally, it presents an exploratory master-key preprocessing construction
based on the trace and norm of the quadratic extension
$\mathrm{GF}(2^{16})/\mathrm{GF}(2^8})$, a permutation polynomial, and a
full-rank linear transformation. The accompanying Python implementation
implements the AES procedure represented by the original project code; the
additional preprocessing construction is discussed mathematically but is not
silently substituted for the standard AES key schedule.
\end{abstract}

\section{Construction of $\mathrm{GF}(256)$}

Let $L$ be a field of characteristic $2$. For every $a\in L$,
\[
a+a=0.
\]
Thus addition has the same XOR behavior as binary addition without carry.

The field $\mathrm{GF}(256)$ can be constructed as
\[
K=\mathrm{GF}(2)[x]/(p(x)),
\]
where
\[
p(x)=x^8+x^4+x^3+x+1.
\]
The polynomial is irreducible over $\mathrm{GF}(2)$, so $(p(x))$ is maximal
and $K$ is a field. A basis is
\[
\{1,x,x^2,x^3,x^4,x^5,x^6,x^7\}.
\]
Consequently,
\[
|K|=2^8=256.
\]

Each byte corresponds to one polynomial in this basis. Multiplication is
ordinary polynomial multiplication over $\mathrm{GF}(2)$ followed by
reduction modulo $p(x)$.

\section{Polynomial Arithmetic and Multiplicative Inversion}

The ring $\mathrm{GF}(2)[x]$ is Euclidean. For
$f(x)$ satisfying
\[
\gcd(f(x),p(x))=1,
\]
the extended Euclidean algorithm gives
\[
u(x)f(x)+v(x)p(x)=1.
\]
Modulo $p(x)$,
\[
u(x)f(x)\equiv1\pmod{p(x)},
\]
so $u(x)$ represents the inverse of $f(x)$ in $K$.

The Python implementation in \texttt{src/aes.py} represents polynomials as
integers whose binary digits are their coefficients. XOR implements addition
over $\mathrm{GF}(2)$.

\section{The AES S-box}

The standard AES S-box is constructed from multiplicative inversion in
$\mathrm{GF}(2^8)$ followed by an affine transformation. For an input
$a\neq0$, first compute $a^{-1}$ in $\mathrm{GF}(2^8)$; for $a=0$, the inverse
is defined as $0$ for the S-box construction. The affine transformation uses
the constant
\[
c=\texttt{0x63}.
\]

The implementation generates all 256 S-box values programmatically rather
than storing a lookup table.

\section{AES-128 Key Expansion}

A 16-byte master key is divided into four words. The AES-128 expansion
generates 44 words, corresponding to 11 round keys. Every fourth word is
processed using \texttt{RotWord}, \texttt{SubWord}, and the appropriate round
constant:
\[
w_i=w_{i-4}\oplus \operatorname{Temp},
\]
where the transformation of $\operatorname{Temp}$ includes the S-box and
the round constant.

The round constants used by the implementation are
\[
01,02,04,08,10,20,40,80,1B,36.
\]

\section{AES Round Transformations}

A 16-byte block is represented as a $4\times4$ state matrix in column-major
order.

\subsection{SubBytes}

Each state byte is replaced using the AES S-box.

\subsection{ShiftRows}

Row $0$ is unchanged, while rows $1$, $2$, and $3$ are shifted left by
$1$, $2$, and $3$ bytes, respectively.

\subsection{MixColumns}

Each column is multiplied over $\mathrm{GF}(256)$ by
\[
M=
\begin{pmatrix}
02&03&01&01\\
01&02&03&01\\
01&01&02&03\\
03&01&01&02
\end{pmatrix}.
\]

The inverse transformation uses
\[
M^{-1}=
\begin{pmatrix}
0E&0B&0D&09\\
09&0E&0B&0D\\
0D&09&0E&0B\\
0B&0D&09&0E
\end{pmatrix}.
\]

\subsection{AddRoundKey}

The state and round key are combined bytewise using XOR.

For AES-128, rounds $1$ through $9$ apply SubBytes, ShiftRows, MixColumns,
and AddRoundKey. The tenth round omits MixColumns.

\section{Padding and Encryption Mode}

The accompanying implementation applies PKCS\#7 padding and encrypts each
16-byte block independently. Thus the implemented mode is ECB.

ECB is retained because it reflects the structure of the original project
code. It is not recommended for production use because repeated plaintext
blocks produce repeated ciphertext blocks.

\section{Proposed Master-Key Preprocessing}

The original mathematical proposal explores the extension
\[
\mathrm{GF}(2^{16})/\mathrm{GF}(2^8).
\]
Let $\theta$ satisfy
\[
\theta^2+\theta=33.
\]
Represent a paired pair of $\mathrm{GF}(2^8)$ elements as
\[
\beta=a+b\theta.
\]

The Frobenius automorphism maps $\theta$ to its conjugate $\theta+1$.
Therefore,
\[
\begin{aligned}
\operatorname{Tr}_{\mathrm{GF}(2^{16})/\mathrm{GF}(2^8)}(\beta)
&=\beta+\beta^{2^8}\\
&=(a+b\theta)+(a+b(\theta+1))\\
&=b.
\end{aligned}
\]
Hence
\[
\boxed{
\operatorname{Tr}_{\mathrm{GF}(2^{16})/\mathrm{GF}(2^8)}(\beta)=b.
}
\]

Similarly,
\[
\begin{aligned}
N_{\mathrm{GF}(2^{16})/\mathrm{GF}(2^8)}(\beta)
&=\beta\beta^{2^8}\\
&=(a+b\theta)(a+b(\theta+1))\\
&=a^2+ab+b^2(\theta^2+\theta)\\
&=a^2+ab+33b^2.
\end{aligned}
\]
Thus
\[
\boxed{
N_{\mathrm{GF}(2^{16})/\mathrm{GF}(2^8)}(\beta)
=a^2+ab+33b^2.
}
\]

The proposed nonlinear transformation is
\[
v_i=u_i^7+t_i u_i+n_i+C_i,
\]
where
\[
t_i=\operatorname{Tr}(u_i),\qquad
n_i=N(u_i).
\]
Since
\[
\gcd(7,2^{16}-1)=1,
\]
the map
\[
u\mapsto u^7
\]
is a permutation of $\mathrm{GF}(2^{16})$. A full-rank matrix
\[
M\in M_8(\mathrm{GF}(256))
\]
can then be used for a linear mixing step
\[
V\mapsto MV.
\]

This construction is an experimental proposal. It is not part of the
standard AES specification, and its cryptographic security would require
separate analysis.

\section{Implementation and Reproducibility}

The Python implementation is organized as follows:
\begin{itemize}
    \item \texttt{poly\_degree}: polynomial degree;
    \item \texttt{F2\_mult}: polynomial multiplication over $\mathrm{GF}(2)$;
    \item \texttt{F2\_divi}: polynomial division with remainder;
    \item \texttt{GF256\_mult}: multiplication in $\mathrm{GF}(256)$;
    \item \texttt{GF256\_inve}: multiplicative inversion;
    \item \texttt{generate\_sbox}: generation of the AES S-box;
    \item \texttt{key\_expansion}: AES-128 key expansion;
    \item \texttt{sub\_bytes}, \texttt{shift\_rows}, \texttt{mix\_columns},
          \texttt{add\_round\_key}: AES round transformations;
    \item \texttt{AES\_encrypt}: block encryption with PKCS\#7 padding;
    \item \texttt{AES\_decrypt}: inverse AES transformation.
\end{itemize}

The example uses
\[
K=\texttt{000102030405060708090a0b0c0d0e0f}
\]
and the plaintext ``mathematical science''. Running
\texttt{run\_example.py} verifies that encryption followed by decryption
recovers the original plaintext.

\section{Discussion and Limitations}

The project is based on undergraduate algebra, finite-field theory, and the
study of AES. The standard AES implementation and the additional mathematical
proposal should be distinguished clearly: the former is an implementation of
the AES procedure represented by the original code, whereas the latter is an
experimental construction for further investigation.

The proposed preprocessing was not fully integrated into the original Python
encryption pipeline. Consequently, this repository does not claim that the
modified construction is equivalent to, or as secure as, standard AES.

\section{AI Assistance}

AI assistance was used for some complex implementation details, including
polynomial division, multiplicative inversion, and linear transformations.
It was also used to identify certain implementation issues. The decryption
procedure was developed primarily through function reuse and derivation of
the corresponding inverse operations.

\begin{thebibliography}{9}

\bibitem{nie}
L. Nie and S. Ding,
\emph{An Introduction to Algebra},
Higher Education Press, Beijing, 2000.

\bibitem{feng}
R. Feng,
\emph{Foundations of Finite Fields},
Higher Education Press, Beijing, 2024.

\end{thebibliography}

\end{document}
